Courts have applied a LW $> 3$ threshold in redistricting challenges.
We survey the documented cases and expert reports.

\subsection{Illinois (7th Circuit)}

Redistricting litigation in Illinois has produced expert reports using LW
to document elongation.
In challenges to Illinois congressional districts following the 2010 Census,
expert witnesses reported LW values for challenged districts using MBR-based
computation.
Districts with LW exceeding 3.0 were characterised as ``elongated'' and
used to support claims that the districts were drawn to connect distant
Democratic-leaning communities along thin corridors~\citep{niemi1990measuring}.

\subsection{New York State Courts}

New York state redistricting litigation following the 2020 Census included
expert reports citing LW as an elongation metric.
The New York Court of Appeals' analysis of legislative redistricting maps
in 2022 cited compactness criteria including elongation measures,
with expert reports documenting LW values for challenged districts.

\subsection{Threshold Interpretation}

The LW $> 3$ threshold has appeared in litigation as follows:
\begin{itemize}
  \item A district with LW $= 3$ is three times as long as it is wide---an
    aspect ratio that is difficult to justify on geographic grounds unless
    the district follows a natural geographic corridor (e.g., a river valley).
  \item LW $> 4$ suggests an extreme elongation that courts have found
    characteristic of ``finger'' or ``tentacle'' districts designed to
    capture distant population clusters.
  \item LW $\leq 2.5$ is generally regarded as not elongated; values between
    $2.5$ and $3.0$ are contested.
\end{itemize}

\subsection{Caveat: No Federal Legal Standard}

No federal court has established LW $> 3$ as a binding legal threshold.
The standard appears in expert reports and state court decisions as a
heuristic, not a categorical rule.
Practitioners should present LW as one component of a multi-metric compactness
profile (K.7) rather than as a pass-fail criterion.
